% Ad Familiares 10.31 (ad-familiares-10-31)
% Source: https://raw.githubusercontent.com/PerseusDL/canonical-latinLit/master/data/phi0474/phi056/phi0474.phi056.perseus-lat1.xml
% Fetched by scripts/fetch_latin.py

% Dateline (Perseus): Scr. Cordubae xvii K. Apr. a. 711 (43).

\ciceroLetterOpener{C. ASINIVS POLLIO CICERONI S. D.}

\ciceroSection{1} minime mirum tibi debet videri nihil me scripsisse de re p., postea quam itum est ad arma. nam saltus Castulonensis, qui semper tenuit nostros tabellarios, etsi nunc frequentioribus latrociniis infestior factus est, tamen nequaquam tanta in hora est, quanta qui locis omnibus dispositi ab utraque parte scrutantur tabellarios et retinent. itaque nisi nave perlatae litterae essent, omnino nescirem quid istic fieret. nunc vero nactus occasionem, postea quam navigari coeptum est, cupidissime et quam creberrime potero scribam ad te.

\ciceroSection{2} ne movear eius sermonibus quem tametsi nemo est qui videre velit, tamen nequaquam proinde ac dignus est oderunt homines, periculum non est ; adeo est enim invisus mihi, ut nihil non acerbum putem, quod commune cum illo sit. natura autem mea et studia trahunt me ad pacis et libertatis cupiditatem. itaque illud initium civilis belli saepe deflevi ; cum vero non liceret mihi nullius partis esse, quia utrubique magnos inimicos habebam, ea castra fugi, in quibus plane tutum me ab insidiis inimici sciebam non futurum ; compulsus eo quo minime volebam, ne in extremis essem, plane pericula non dubitanter adii.

\ciceroSection{3} Caesarem vero, quod me in tanta fortuna modo cognitum vetustissimorum familiarium loco habuit, dilexi summa cum pietate et fide. quae mea sententia gerere mihi licuit, ita feci ut optimus quisque maxime probarit; quod iussus sum, eo tempore atque ita feci ut appareret invito imperatum esse. cuius facti iniustissima invidia erudire me potuit quam iucunda libertas et quam misera sub dominatione vita esset. ita si id agitur ut rursus in potestate omnia unius sint, quicumque is est, ei me profiteor inimicum, nec periculum est ullum quod pro libertate aut refugiam aut deprecer.

\ciceroSection{4} sed consules neque senatus consulto neque litteris suis praeceperant mihi quid facerem ; unas enim post Idus Mart. demum a Pansa litteras accepi, in quibus hortatur me ut senatu scribam me et exercitum in potestate eius futurum. quod , cum Lepidus contionaretur atque omnibus scriberet se consentire cum Antonio, maxime contrarium fuit ; nam quibus commeatibus invito illo per illius provinciam legiones ducerem? aut si cetera transissem, num etiam Alpis poteram transvolare, quae praesidio illius tenentur? adde huc quod perferri litterae nulla condicione potuerunt ; sescentis enim locis excutiuntur, deinde etiam retinentur ab Lepido tabellarii.

\ciceroSection{5} illud me Cordubae pro contione dixisse nemo vocabit in dubium, provinciam me nulli, nisi qui ab senatu missus venisset, traditurum. nam de legione tricesima tradenda quantas contentiones habuerim quid ego scribam? qua tradita quanto pro re p. infirmior futurus fuerim quis ignorat? hac enim legione noli acrius aut pugnacius quicquam putare esse. qua re eum me existima esse, qui primum pacis cupidissimus sim (omnis enim civis plane studeo esse salvos), deinde qui et me et rem publicam vindicare in libertatem paratus sim.

\ciceroSection{6} quod familiarem meum tuorum numero habes, opinione tua mihi gratius est ; invideo illi tamen, quod ambulat et iocatur tecum. quaeres quanti aestimem. si umquam licuerit vivere in otio, experieris ; nullum enim vestigium abs te discessurus sum. illud vehementer admiror, non scripsisse te mihi manendo in provincia an ducendo exercitum in Italiam rei p. magis satis facere possim. ego quidem, etsi mihi tutius ac minus laboriosum est manere, tamen, quia video tali tempore multo magis legionibus opus esse quam provinciis, quae praesertim reciperari nullo negotio possunt, constitui, ut nunc est, cum exercitu proficisci. deinde ex litteris, quas ' Pansae misi, cognosces omnia ; nam tibi earum exemplar misi. xvii K. Apr. Corduba.
